% Draft of the 2026-10 rewrite: Introduction and Related Work. Numbers marked \pending{} wait for the round-1/2 arms.
% To be merged into main.tex once the method results are in. Measurement numbers come from fair_cdfb.json,
% cdf_filtered.json and dfd_filtered.json (one rule for every detector).

\begin{abstract}
Many of the faces a forgery detector receives have been beautified by a camera app or a retouching service. We measure
what this does to thirteen published detectors with one rule for all of them: the score is the probability of forgery,
the threshold is set on untreated genuine faces, and both errors are counted, genuine faces accused and forgeries let
through. On three sources that none of the detectors was tuned on, app presets on Celeb-DF-B, eleven colour filters on
Celeb-DF-v2 and sixteen further filters on DFD, every detector loses AUC under beautification; for the two CLIP-based
detectors of 2025 the miss rate at a fixed threshold rises by 7 to 17 points, and a retouching service makes eleven
detectors accuse 36--63\% of its most strongly smoothed genuine faces. A probe of the feature space shows where the damage
comes from: a radial luminance overlay moves a genuine face along the forgery direction three to ten times further than a
pure colour look, and the direction cannot be projected out without losing the forgery evidence itself. We therefore train
against it. We define a bounded, differentiable operator over identity-preserving edits whose range is fitted to measured
retouching and filters, and minimise the risk under the worst edit the operator can produce for each face, found by
gradient ascent, on genuine and forged faces alike. With FaceForensics++ as the only training data, \pending{the worst-case
model reaches an AUC of X under app presets and Y under unseen filters, against 0.81/0.79 for Effort and 0.83/0.79 for
Forensics Adapter}, while matching them on untreated Celeb-DF-v2 and DFD. Random edits from the same operator and the worst
of four random draws fall short of the gradient search, and a model that also learns from commercial retouching pairs
extends the margin. The detector's evidence map names regions whose restoration removes \pending{Z}\% of its detections on
held-out part-level forgeries, within a point of the true footprint.
\end{abstract}

% ----------------------------------------------------------------------------
\section{Introduction}
Face-forgery detectors are used for content moderation, remote identity checks and fact-checking. In these settings a
genuine face called fake accuses a real person, and a forgery called real is let through. Many of the genuine faces such a
detector receives have been beautified. Retouching is common on social media and photo-sharing sites and appears even on
identification documents~\cite{bharati2016}; skin smoothing, whitening, eye enlargement and face slimming are default
options in camera apps and are sold by retouching APIs, and colour presets ship with every photo-sharing app.
Beautification edits the pixels of the face, which is what a forgery detector inspects, but changes neither who is shown
nor what the person does.

We measure both errors on thirteen published detectors, from Xception to the CLIP-based Effort~\cite{effort} and Forensics
Adapter~\cite{forensicsadapter}, with one rule for every detector: the score is its probability of forgery, the threshold is
the one that accuses 5\% of untreated genuine frames, and a forgery not called fake is a miss whatever else it was called.
Three test sources carry beautification that none of the detectors was trained on. Two of them, Celeb-DF-B and the DFD
filters, were never used for any design decision of ours; the Celeb-DF-v2 filters informed the design of our edit operator
and are reported as a development source. On Celeb-DF-B~\cite{libourel2024},
whose app presets were applied by its authors, re-encoding alone explains part of the change, so we subtract it; the miss
rate still rises by 11.0, 15.2 and 21.8 points for Effort, Forensics Adapter and SBI~\cite{sbi}. On Celeb-DF-v2 frames under
eleven colour filters never used in our training, the AUC of the two CLIP detectors falls from 0.88--0.89 to 0.79. On DFD
frames under sixteen further filters, built after every design decision of this paper had been taken, it falls from 0.94 to
0.90--0.91 and Effort's miss rate rises by 12 points. The other direction is as real: at a threshold that accuses 5\% of a
retouching service's untouched originals, the eleven FF++-trained detectors accuse 36--63\% of its most strongly smoothed
genuine faces.

Where does the damage come from? We look at the pooled feature of a CLIP detector trained on FF++ only. Of eleven held-out
filters, the two that add a radial luminance overlay move a genuine face's forgery logit by +5.6 and +7.8 on average, the
nine pure colour looks by +0.1 to +2.4, and the overlay shifts have a cosine of 0.20--0.33 with the forgery weight vector
against 0.00--0.09 for the colour looks. A luminance ramp around the face reads as a blending boundary. Most of the shift
energy of every held-out filter lies in the subspace spanned by the shifts of beautifications seen in training, so the
directions were seen and the combinations were not. Projecting that subspace out and refitting the head does not help:
the clean AUC drops at once and the false accusations under the filters do not fall, because the same few directions
carry the forgery evidence. The separation has to be learnt.

We learn it by training against the worst beautification rather than a random one. The space of identity-preserving edits
is written as a differentiable operator $B(x;u)$ with bounded tone curves, a colour matrix, saturation down to monochrome,
contrast, a smooth luminance and tint field, skin smoothing and grain. Its range is not set by hand: the operator is fitted
to the renders of two commercial retouching services and to the colour filters used in training, and the search is
confined to the range of the fitted edits. The training objective is the risk under the worst edit in that range,
\begin{equation*}
\min_\theta\ \mathbb{E}_{(x,y)}\Big[\ell\big(f_\theta(x),y\big)+\max_{u\in U}\ell\big(f_\theta(B(x;u)),y\big)\Big],
\end{equation*}
with the inner maximisation approximated by signed-gradient ascent from a measured edit, applied to genuine and forged
faces alike so that no edit is associated with either label. This is adversarial training with a perturbation set that
has a meaning: not an $\ell_p$ ball but the edits a user can apply. Random edits from the same operator, and the worst of
four random draws, are the controls that separate the contribution of the search from that of the operator family.

Our contributions are:
\begin{itemize}
\item \textbf{A two-sided measurement with one rule.} Thirteen detectors on three beautification sources and a commercial
retouching service, scored by the same rule, with video-level confidence intervals and re-encoding separated from
beautification.
\item \textbf{The mechanism, measured.} A feature-space probe that quantifies which edits move a CLIP detector toward
``forged'', by how much, and why a linear fix cannot work.
\item \textbf{Worst-case beautification training.} A bounded differentiable edit operator fitted to measured retouching and
filters, and robust training against its worst edit on both classes, with the random and worst-of-$k$ controls that
isolate the effect of the search. \pending{With FF++ as the only training data it reaches AUC X / Y / Z under the three
sources against 0.81 / 0.79 / 0.90 for Effort}, at unchanged accuracy on untreated Celeb-DF-v2 and DFD.
\item \textbf{Evidence with a faithfulness test.} An evidence head trained on exact edit footprints, checked by restoring
the region it names and watching the decision change.
\end{itemize}
We also report what did not work, because it bounds the claims. A third output class for beautified faces, which we
proposed in an earlier version of this work, does not improve detection: a detector with the same data and a single forged
output is better on every source by about 0.02 AUC, and the class is kept only as a user-facing output. In our runs, beautification
applied with unequal frequency to the two classes taught the detector a correlation and moved its errors from one side to
the other rather than removing them. And the model still separates lightly retouched faces from their originals at only
70--75\% balanced accuracy on an unseen service.

% ----------------------------------------------------------------------------
\section{Related Work}
\textbf{Beautification and forgery detection.} Libourel et al.~\cite{libourel2024} released Celeb-DF-B and measured AUC
drops of 0.15--0.16 for three detectors under app presets; Saeed et al.~\cite{saeed2025} show that face enhancement lets
deepfakes evade six detectors and fine-tune on fixed enhanced samples; Concas et al.~\cite{concas2025} report both errors
for two CNNs under a smoothing filter. None trains against the edit, and none scores detectors by a common rule. The
retouching literature labels commercial operations~\cite{retouchingffhq} or restores the untouched face~\cite{frrffusion,mofrr,reface,wang2019photoshop}
without a forgery class. Liu et al.~\cite{miea2026} attribute both forgery and beautification tools in a closed set.

\textbf{Attacks on forgery detectors through natural appearance.} StatAttack~\cite{statattack2023} optimises exposure, blur
and noise, and AdvShadow~\cite{advshadow2024} a shadow overlay, to make forgeries pass; both are attacks in one direction
and are not used for training. Fatima et al.~\cite{fatima2025} train on the worst of $K$ random counter-forensic
transforms, chosen by loss without a gradient. We search the edit space by gradient, on both classes, and measure both
errors; the worst-of-$k$ choice is one of our controls.

\textbf{Adversarial training over semantic perturbations.} Hosseini and Poovendran~\cite{hosseini2018}, ColorFool~\cite{colorfool2020},
functional attacks~\cite{laidlaw2019} and colour-filter attacks~\cite{ace2020,advcf2023} move images along hue, saturation
or a parametric colour curve; the vignetting attack~\cite{ava2021} does so with a radial field. Hsiung et al.~\cite{gat2023}
train against compositions of hue, saturation, brightness, contrast and rotation. Adversarial data augmentation for domain
generalisation~\cite{volpi2018} and group-robust optimisation~\cite{sagawa2020} share the objective. Our perturbation set is
the space of beautification edits, its range is fitted to measured edits, and the task is forgery detection, where the
perturbation also appears on the positive class. Single-step adversarial training can overfit to its own step~\cite{wong2020fast};
we test for this with a ten-step attack.

\textbf{Generalisable forgery detection.} Face X-ray~\cite{facexray} and SBI~\cite{sbi} synthesise blending artefacts from
real images, UCF~\cite{ucf} separates common from method-specific forgery features, and Effort~\cite{effort} and Forensics
Adapter~\cite{forensicsadapter} adapt CLIP ViT-L/14 with few trainable parameters. DeepfakeBench~\cite{deepfakebench}
provides common checkpoints. The alpha-blending hypothesis~\cite{alphablend2026} holds that frame-based detectors localise
compositing artefacts; our measurement that a luminance ramp around the face raises the forgery score is consistent with
it, for genuine faces.

\textbf{Explanations for detectors.} GAIN~\cite{gain} trains attention so that erasing the attended region removes the
evidence, amortised maskers predict saliency in one pass~\cite{dabkowski2017,phang2020}, and DOLOS~\cite{dolos} compares
methods for localising inpainted regions. We supervise the map with exact footprints and test it by restoring pixels.
